% Canonical statement: sections/06_explicit.tex:28-38
\begin{theorem}[Conditional almost-every-key explicit lower bounds]\label{thm:prf}
Assume the stated PRF security for a fixed $\delta>0$. Fix $c\ge1$, $0<\epsilon<1/4$, $n=3b+1$, and put
\[
 \sigma=n^{\lceil(4c+2)/\delta\rceil}.
\]
For every key, entries of $A_s$ and hypotheses in $H_{A_s}$ are evaluable in $\operatorname{poly}_c(n)$ time. Every key has RECTIFY and proper learners with $O_\epsilon(n)$ queries and tolerance $\Omega_\epsilon(1)$. For all sufficiently large admissible $n$, outside a negligible fraction of keys,
\[
 \boxed{\dc(H_{A_s})\ge\frac{n^c}{30\log_2n}.}
\]
Here the learner with $O(n)$ queries may use table-polynomial computation. A succinct polynomial-time implementation with $O(n^2)$ queries is established in Section~\ref{sec:fast}.
\end{theorem}

% Proof binding: appendices/explicit_proofs.tex:23-37
\begin{proof}[Proof of \cref{thm:prf}]
The evaluation claim follows from integer arithmetic and one PRF evaluation. The rectangle and proper-learning bounds hold for every monotone flip, without any randomness assumption on its entries.

Take the subgrid scale $k=\lceil c n^c\rceil$, which is at most $N$ for sufficiently large $n$. Restrict both the full row set and point set to $[k]\times[2k^2]$. The template and flip law restricted there are exactly the scale-$k$ construction. An oracle distinguisher reads its incidence bits and accepts precisely when its sign-rank is at most $D_k$. Its time is
\[
 2^{O_c(n^{4c})}\le2^{n^{4c+2}}\le2^{\sigma^\delta}
\]
for sufficiently large $n$. Under a uniform function its acceptance probability is at most $2^{-k^4/2}$ by Lemma~\ref{lem:subgrid}. Under the PRF it is therefore at most $2^{-k^4/2}+\operatorname{negl}(\sigma)$, which is negligible in $n$. Restricting a sign realization cannot increase its rank. Consequently, except on that event, the full dimension is at least $D_k+1>k/(24\log_2k)$.

For large $n$, $k\ge c n^c$ and $\log_2k\le c\log_2n+\log_2(2c)\le(5/4)c\log_2n$. Hence
\[
 D_k+1>\frac{k}{24\log_2k}\ge\frac{n^c}{30\log_2n}.
\]
This proves the displayed constant. The choice $k=n^c$ alone would instead give $n^c/(24c\log_2n)$; its factor $c$ cannot be discarded. Also $N=2^{(n-1)/3}$, rather than $2^{n/3}$, respects the fixed definition $n=\log_2|X|$.
\end{proof}
